\chapter{Introduction}\label{I-introduction}

Le but de ce cours est de consolider et compléter les connaissances en
théorie des probabilités acquises en CPGE mais surtout de permettre
d'acquérir le raisonnement probabiliste. En effet, les probabilités
peuvent être vues comme un outil de modélisation de phénomènes qui ont
la caractéristique d'être aléatoires. L'aléatoire peut intervenir de
différentes manières dans ces phénomènes.

\begin{itemize}
\tightlist
\item
  Dans les cas d'école que sont les jeux de pile ou face ou de lancés de
  dés, la différence entre les résultats, si l'on réitère l'expérience,
  peut être liée à l'impulsion initiale communiquée au dé et à d'autres
  facteurs environnementaux comme le vent, la rugosité de la table, etc.
  Le hasard intervient du fait de la méconnaissance des conditions
  initiales, car la pièce ou le dé ont des trajectoires parfaitement
  définies par la mécanique classique.
\item
  Dans beaucoup de cas de figure, on fait intervenir l'aléatoire dans la
  modélisation du fait d'une connaissance incomplète des phénomènes. On
  parle alors de modélisation de l'incertitude. C'est le cas par exemple
  en sciences du climat, dont nous discutons ci-dessous.
\item
  Dans certains domaines, tels la physique quantique, l'aléatoire fait
  intrinsèquement partie de la théorie.
\end{itemize}

Les probabilités sont aussi un préalable indispensable pour aborder
l'analyse statistique des données et les méthodes d'apprentissage
automatique.

En CPGE, les probabilités ont été vues dans le cadre de phénomènes
aléatoires qui admettent un nombre au plus dénombrable de résultats
possibles. Ce cadre restreint est supposé connu. On pourra se reporter
au chapitre 3 du cours de Garnier
(\protect\hyperlink{refI-polygarnier}{2018}) ou aux deux premiers
chapitres du cours de Jourdain (\protect\hyperlink{refI-polyponts}{2016})
pour une éventuelle mise à niveau.

\section{Exemple de la modélisation du
climat}\label{I-exemple-de-la-moduxe9lisation-du-climat}

Les prédictions de la météo et les projections du climat proviennent
généralement de modèle numériques qui simulent les différents processus
physiques à l'oeuvre. Les incertitudes dans la construction et
l'application de ces modèles sont variées et peuvent être réparties en
quatre groupes~: les conditions initiales (on ne connaît jamais
parfaitement l'ensemble des variables climatiques en tout point du
globe), les conditions aux limites (par exemple lorsqu'on travaille à
l'échelle d'un continent ou d'un pays), les valeurs des paramètres
intervenant dans les modèles (constantes issues d'observations
diverses), les incertitudes structurelles enfin qui relèvent des choix
de modélisation. Pour tenir compte de ces incertitudes, les
climatologues effectuent des ensembles de simulations, où les
différentes quantités incertaines sont échantillonnées selon des modèles
probabilistes (voir en particulier Tebaldi and Knutti
(\protect\hyperlink{refI-tebaldi}{2007})).

L'incertitude sur les conditions initiales est particulièrement
influente aux faibles échelles de temps. La météo est un système
chaotique, les prévisions sont extrêmement sensibles aux variations des
conditions initiales utilisées pour initialiser les modèles. Ces
dernières sont en revanche connues de manière imparfaite. Il est ainsi
nécessaire de tenir compte de cette incertitude. Ceci est rendu possible
par la modélisation probabiliste. Les conditions aux limites font
intervenir des \textbf{variables continues} (température, pression,
vitesse du vent, etc.), ce qui représente la principale nouveauté par
rapport au programme de CPGE, et il est nécessaire de caractériser leurs
relations de \textbf{dépendance}. Il convient par ailleurs de tenir
compte des observations (satellites, station de mesure, \ldots) en
incluant cette information dans la modélisation probabiliste. On parle
de \textbf{conditionnement} aux données. Il s'agit enfin d'en générer
les valeurs via des algorithmes de \textbf{simulation stochastique}. La
validité de l'approche est assurée par les \textbf{théorèmes limites},
qui garantissent la représentativité des ensembles générés.

\section{Plan du cours}\label{I-plan-du-cours}

Le cours est organisé en 5 amphis et abordera consécutivement les
notions évoquées ci-dessus, à savoir les probabilités définies sur
\(\mathbb{R}\), les variables et vecteurs aléatoires réels,
l'indépendance et le conditionnement de variables aléatoires, l'étude
des suites de variables aléatoires et, enfin, les méthodes de simulation
stochastique.

\section{Historique}\label{I-historique}

Avant que l'étude des probabilités soit considérée comme une science,
l'observation du hasard dans les événements naturels a amené les
philosophes et les scientifiques à réfléchir sur la notion de liens
entre événements, causes et conséquences, et lois de la nature. Les jeux
de hasard, les situations météorologiques ou les trajectoires des astres
ont fait partie des domaines étudiés. Les explications données sont
alors liées au destin, à une colère céleste ou à une présence divine.

Il est communément admis que le début de la science des probabilités se
situe au XVIe siècle avec l'analyse de jeux de hasard par Jérôme Cardan
et au XVIIe siècle avec les discussions entre Pierre de Fermat et Blaise
Pascal au sujet de paradoxes issus de ces jeux, notamment posés par
Antoine Gombaud, chevalier de Méré. Cette nouvelle théorie est nommée
géométrie aléatoire par le chevalier de Méré en 1654, elle est appelée
par la suite calcul conjectural, arithmétique politique et plus
communément aujourd'hui théorie des probabilités. Cette théorie, dite
des probabilités modernes, est alors étudiée par de nombreux penseurs
jusqu'au XIXe siècle~: Kepler, Galilée, Leibniz, Huygens, Halley,
Buffon, les frères Bernoulli, Moivre, Euler, D'Alembert, Condorcet,
Laplace, Fourier. Elle est principalement basée sur les événements
discrets et la combinatoire.

Des considérations analytiques ont forcé l'introduction de variables
aléatoires continues dans la théorie. Cette idée prend tout son essor
dans la théorie moderne des probabilités, dont les fondations ont été
posées par Andreï Nikolaevich Kolmogorov. Kolmogorov combina la notion
d'univers, introduite par Richard von Mises et la théorie de la mesure
pour présenter son système d'axiomes pour la théorie des probabilités en
1933. Très vite, son approche devint la base incontestée des
probabilités modernes.

Le XXe siècle voit également le développement de l'application de la
théorie des probabilités dans plusieurs sciences.

Avec la mécanique newtonienne, la théorie du champ électromagnétique ou
la thermodynamique, la physique classique est la théorie utilisée
jusqu'à la fin du XIXe siècle. En 1925, Erwin Schrödinger étudie
l'équation qui détermine l'évolution d'une onde au cours du temps~:
l'équation de Schrödinger. Max Born utilise cette équation pour décrire
une collision entre des particules telles que des électrons ou des
atomes. Les observations de ces expériences l'amènent à supposer que le
module de la fonction d'onde est la probabilité que la particule soit
détectée en un point de l'espace. C'est le début d'une nouvelle approche
de la physique quantique.

En 1900, Louis Bachelier fut un des premiers mathématiciens à modéliser
les variations de prix boursiers grâce à des variables aléatoires. « le
marché n'obéit qu'à une seule loi~: la loi du hasard ». Bachelier
utilise alors le calcul stochastique pour étudier les variations
boursières au cours du temps. En 1970, Fischer Black et Myron Scholes
reprennent les idées de Bachelier pour modéliser les rendements d'une
action.

L'utilisation des probabilités en biologie a pris un essor dans les
années 1970, notamment dans l'étude de l'évolution des espèces. La
reproduction des individus est modélisée par un choix aléatoire des
gènes transmis ainsi que des mutations apparaissant de manière aléatoire
sur les individus. L'extinction des espèces ou des gènes est alors
étudiée en fonction des effets stochastiques.

De nos jours, l'Ecole française de Probabilités est très active. La
première Médaille Fields décernée à un probabiliste a été attribuée à
Wendelin Werner en 2006. Les probabilités se développent de plus en
plus, alimentées en particulier de manière essentielle par la physique,
le développement des réseaux de télécommunications, la finance, la
biologie, la médecine\ldots{} Elles permettent de construire des modèles
mathématiques, qui peuvent être validés par les données suivant la
théorie statistique, et fournissent également des possibilités
d'expérimentations fictives dans de multiples domaines d'applications.

\chapter{Objectifs d'apprentissage}\label{I-objectifs-dapprentissage}

\runinhead{Probabilités des
événements}\phantomsection\label{I-probabilituxe9s-des-uxe9vuxe9nements}

\begin{itemize}
\item
  connaître les définitions de base des probabilités~:

  \begin{itemize}
  \tightlist
  \item
    $\mathord{\boldsymbol{\circ}}$ espace fondamental \(\Omega\)
  \item
    $\mathord{\boldsymbol{\circ}}$ événement \(A \subset \Omega\)
  \item
    $\mathord{\boldsymbol{\circ}}$ tribu (de parties)
  \item
    $\mathord{\bullet}$ probabilité
  \item
    $\mathord{\bullet}$ propriété presque sûre
  \end{itemize}
\item
  connaître, savoir démontrer et exploiter les résultats suivants~:

  \begin{itemize}
  \tightlist
  \item
    $\mathord{\bullet}$ les propriétés élémentaires d'une probabilité
  \item
    $\mathord{\bullet}$ le théorème de continuité monotone
  \end{itemize}
\item
  indépendance et conditionnement

  \begin{itemize}
  \tightlist
  \item
    $\mathord{\bullet}$ connaître la définition de la probabilité
    conditionnelle
  \item
    $\mathord{\bullet}$ connaître et savoir exploiter la formule des
    probabilités totales
  \item
    $\mathord{\bullet}$ connaître et savoir exploiter la formule de
    Bayes
  \item
    $\mathord{\bullet}$ connaître la définition de l'indépendance
    d'événements
  \end{itemize}
\end{itemize}

\runinhead{Variables
aléatoires}\phantomsection\label{I-variables-aluxe9atoires}

\begin{itemize}
\item
  variables aléatoires réelles

  \begin{itemize}
  \tightlist
  \item
    $\mathord{\bullet}\mathord{\bullet}$ connaître la définition de la
    tribu borélienne
  \item
    $\mathord{\bullet}\mathord{\bullet}$ savoir manipuler la notion de
    tribu engendrée
  \item
    $\mathord{\bullet}\mathord{\bullet}$ savoir que la tribu borélienne
    est engendrée par les intervalles de la forme
    \(]-\infty, x],~ x\in \mathbb{R}\)
  \item
    $\mathord{\bullet}$ connaître la définition d'une variable aléatoire
    réelle et de sa loi de probabilité
  \end{itemize}
\item
  loi des variables aléatoires réelles

  \begin{itemize}
  \tightlist
  \item
    $\mathord{\bullet}$ savoir que la fonction de répartition
    caractérise la loi d'une variable aléatoire réelle
  \item
    $\mathord{\bullet}\mathord{\bullet}$ connaître les propriétes de
    base de la fonction de répartition
  \item
    $\mathord{\bullet}\mathord{\bullet}$ savoir calculer la fonction de
    répartition d'une variable aléatoire
  \end{itemize}
\item
  densité de probabilité

  \begin{itemize}
  \tightlist
  \item
    $\mathord{\bullet}\mathord{\bullet}$ connaître la définition d'une
    densité de probabilité
  \item
    $\mathord{\bullet}\mathord{\bullet}$ connaître la définition d'une
    variable aléatoire réelle à densité
  \item
    $\mathord{\bullet}\mathord{\bullet}$ savoir identifier si une
    variable aléatoire admet une densité
  \end{itemize}
\end{itemize}

\chapter{Probabilités des
événements}\label{I-probabilituxe9s-des-uxe9vuxe9nements-1}

\section{Phénomènes aléatoires et
événements}\label{I-phuxe9nomuxe8nes-aluxe9atoires-et-uxe9vuxe9nements}

L'objet de la théorie des probabilités est l'analyse mathématique de
phénomènes dans lesquels le hasard intervient. Les phénomènes aléatoires
résultent d'expériences dont le résultat ne peut être prédit à l'avance
et qui \textbf{peut} varier si on répète l'expérience dans des
conditions identiques.

Il est aisé de trouver des exemples de tels phénomènes.

\begin{exemple}\phantomsection\label{I-exemples}
\begin{enumerate}
\def\labelenumi{\arabic{enumi}.}
\tightlist
\item
  Jeu de Pile ou Face
\item
  Lancé de dés
\item
  Durée de vie d'une ampoule électrique
\item
  Température demain à 12h au sommet de la tour Eiffel
\item
  Évolution de la vitesse d'une molécule dans un gaz raréfié sur un
  intervalle de temps \([t_1,t_2]\)
\end{enumerate}

La théorie des probabilités vise à fournir un modèle mathématique pour
décrire ces phénomènes. Elle repose sur trois ingrédients essentiels
dont on donne ici les définitions.
\end{exemple}

\addcontentsline{toc}{subsection}{L'espace fondamental}
\begin{definition}[{L'espace
fondamental}]\phantomsection\label{I-lespace-fondamental}

Noté habituellement \(\Omega\), l'\emph{espace fondamental} (ou encore
l'\emph{espace d'état} ou \emph{univers}) contient l'ensemble de tous
les résultats possibles d'un phénomène aléatoire. Un résultat possible
d'une expérience sera noté \(\omega \in \Omega\).

Si on reprend les exemples précédents, on peut facilement définir les
univers associés.
\end{definition}

\begin{exemple}\phantomsection\label{I-exemples-1}
\begin{enumerate}
\def\labelenumi{\arabic{enumi}.}
\tightlist
\item
  Jeu de Pile ou Face, \(\Omega = \{\text{pile}, \text{face}\}\)
\item
  Lancé de dés, \(\Omega = \{1,2,3,4,5,6\}\)
\item
  Durée de vie d'une ampoule électrique, \(\Omega = [0,+\infty [\)
\item
  Température demain à 12h au sommet de la tour Eiffel (en degrés
  Kelvin), \(\Omega = [0,+\infty [\)
\item
  Évolution de la vitesse d'une molécule dans un gaz raréfié sur un
  intervalle de temps \([t_1,t_2]\), \(\Omega\)~: ensemble des
  application continues sur \([t_1,t_2]\) à valeurs dans
  \(\mathbb{R}^3\)
\end{enumerate}

Cette liste d'exemples montre que l'espace \(\Omega\) peut varier
énormément dans sa structure, d'une expérience à l'autre. Cela permet de
réaliser la richesse de la théorie qu'il faut mettre en place, pour
créer un modèle qui englobe tous ces cas. Nous verrons également
ultérieurement que le modèle abstrait que nous allons construire
permettra de s'affranchir du fait que \(\Omega\) décrit précisément tous
les résultats possibles de l'expérience.
\end{exemple}

\addcontentsline{toc}{subsection}{Evénement}
\begin{definition}[{Evénement}]\phantomsection\label{I-evuxe9nement}

Un \emph{événement} est une propriété qui est vérifiée ou non une fois
l'expérience réalisée. On identifie un événement \(A\) à un
sous-ensemble ou \emph{partie} de \(\Omega\),
i.e.~\(A = \{\omega \in \Omega : A \text{ est vérifiée pour } \omega \}\).
\end{definition}

\begin{exemple}\phantomsection\label{I-exemples-2}
\begin{enumerate}
\def\labelenumi{\arabic{enumi}.}
\tightlist
\item
  Jeu de Pile ou Face~: \(A = \{\text{pile}\}\).
\item
  Lancé de dés~: \(A = \{1,3,5\}\).
\item
  Durée de vie d'une ampoule électrique~:
  \(A = [t_1,t_2] \subset \mathbb{R}_+\).
\item
  Température demain à 12h au sommet de la tour Eiffel (en degrés
  Kelvin)~: \(A = [T_1,T_2]\cup[T_3,T_4] \subset \mathbb{R}_+\).
\item
  Évolution de la vitesse d'une molécule dans un gaz raréfié sur un
  intervalle de temps \([t_1,t_2] \subset \mathbb{R}_+\)~:
  \(A = \{ f \in C([t_1,t_2], \mathbb{R}^3) : \|f-g\|_{\infty} \leq a \}\),
  où \(g \in C([t_1,t_2],\mathbb{R}^3)\) et \(a \in \mathbb{R}_+\).
\end{enumerate}

Les événements étant des ensembles, les opérations ensemblistes
classiques admettent une interprétation probabiliste.
\end{exemple}

\subsection*{Correspondance entre opérations logiques et
ensemblistes}\label{I-correspondance-entre-opuxe9rations-logiques-et-ensemblistes}

\begin{longtable}[]{@{}ccc@{}}
\toprule
Terminologie probabiliste & Terminologie ensembliste &
Notation\tabularnewline
\midrule
\endhead
événement certain & ensemble entier & \(\Omega\)\tabularnewline
événement impossible & ensemble vide & \(\varnothing\)\tabularnewline
événement contraire & complémentaire & \(A^c\)\tabularnewline
événement atomique & singleton & \(\{\omega\}\)\tabularnewline
implication & inclusion & \(\subset\)\tabularnewline
et & intersection & \(\cap\)\tabularnewline
ou & réunion & \(\cup\)\tabularnewline
événements incompatibles & ensembles disjoints &
\(A_1\cap A_2 = \varnothing\)\tabularnewline
\bottomrule
\end{longtable}

On doit maintenant répondre à la question de savoir quels sont les
événements dont on va vouloir évaluer la probabilité d'occurence. On va
ainsi regrouper les événements en un ensemble \(\mathcal{A}\) qui
constitue une collection de sous-ensembles de \(\Omega\). On va
souhaiter en particulier pouvoir combiner des événements au sein de
\(\mathcal{A}\) par les opérations ensemblistes courantes.

Ceci conduit à la notion de \emph{tribu de parties} de \(\Omega\).

\addcontentsline{toc}{subsection}{Tribu}
\begin{definition}[{Tribu}]\phantomsection\label{I-deftribu}

Une \emph{tribu} (ou \emph{\(\sigma\)-algèbre}) \(\mathcal{A}\) est une
collection de sous-ensembles de \(\Omega\) tels que~:

\begin{enumerate}
\def\labelenumi{\arabic{enumi}.}
\tightlist
\item
  \(\Omega \in \mathcal{A}\),
\item
  \(A \in \mathcal{A}\Rightarrow A^c \in \mathcal{A}\),
\item
  \(\forall\, n \in \mathbb{N}, A_n \in \mathcal{A}\Rightarrow \bigcup_n A_n \in \mathcal{A}\).
\end{enumerate}
\end{definition}

Le couple \((\Omega, \mathcal{A})\) est appelé \emph{espace
probabilisable}. On notera qu'il s'agit d'un espace mesurable dans la
terminologie du cours d'intégration.

\begin{exemple}\phantomsection\label{I-exemples-3}
\begin{enumerate}
\def\labelenumi{\arabic{enumi}.}
\tightlist
\item
  \(\mathcal{A}= \{\varnothing,\Omega\}\) est la tribu \emph{grossière}
  ou \emph{triviale}~: c'est la plus petite tribu de \(\Omega\).
\item
  Dans le cas où \(\Omega\) est au plus dénombrable, on choisit
  systématiquement l'ensemble \(\mathcal{P}(\Omega)\) des parties de
  \(\Omega\) dont on vérifie aisément qu'il s'agit d'une tribu\,; c'est
  le cadre étudié en CPGE. On verra plus loin que cette tribu est trop
  grande dans le cas où \(\Omega\) est infini non dénombrable.
\item
  Si \(\Omega = \mathbb{R}\), on peut le munir de la tribu formée des
  ensembles mesurables de \(\mathbb{R}\), dite \emph{tribu de Lebesgue}.
\item
  On verra plus loin qu'il est aisé de définir une tribu sur tout espace
  topologique.
\end{enumerate}
\end{exemple}

\section{Probabilité}\label{I-probabilituxe9}

\addcontentsline{toc}{subsection}{Probabilité}
\begin{definition}[{Probabilité}]\phantomsection\label{I-defproba}

Une \emph{probabilité} sur l'espace \((\Omega, \mathcal{A})\) est une
application \(\mathbb{P}: \mathcal{A}\rightarrow [0,1]\), telle que~:

\begin{enumerate}
\def\labelenumi{\arabic{enumi}.}
\tightlist
\item
  \(\mathbb{P}(\Omega) = 1\),
\item
  pour toute suite (dénombrable) \((A_n)_{n\in\mathbb{N}^\star}\)
  d'éléments de \(\mathcal{A}\) \textbf{deux à deux disjoints}, on a
  \begin{equation*}
   \mathbb{P}\left(\bigcup_{n\in\mathbb{N}^\star} A_n\right) = \sum_{n\in\mathbb{N}^\star} \mathbb{P}(A_n).
   \end{equation*}
\end{enumerate}
\end{definition}

Le triplet \((\Omega, \mathcal{A}, \mathbb{P})\) est appelé \emph{espace
probabilisé}. La modélisation probabiliste consiste ainsi à décrire une
expérience aléatoire par la donnée d'un espace probabilisé.

On notera que dans la terminologie du cours d'intégration,
\((\Omega, \mathcal{A}, \mathbb{P})\) est un espace mesuré. On vérifie
aisément que \(\mathbb{P}\) est une mesure qui a la particularité d'être
finie, c'est-à-dire telle que \(\mathbb{P}(\Omega) =1\). On parle aussi
de \emph{mesure de probabilité}.

La définition suivante est fondamentale en théorie des probabilités.
Elle introduit une notion de ``vrai ou faux'' qui dépend de la
probabilité choisie sur l'espace fondamental.

\addcontentsline{toc}{subsection}{Propriété presque-sûre}
\begin{definition}[{Propriété
presque-sûre}]\phantomsection\label{I-propriuxe9tuxe9-presque-suxfbre}

Soit \((\Omega, \mathcal{A}, \mathbb{P})\) un espace probabilisé. On dit
qu'un événement \(A\in\mathcal{A}\) se réalise
\emph{\(\mathbb{P}\)-presque sûrement} (en abrégé \(\mathbb{P}\)-p.s.)
si \(\mathbb{P}(A) = 1\).
\end{definition}

\addcontentsline{toc}{subsection}{Propriétés élémentaires}
\begin{proposition}[{Propriétés
élémentaires}]\phantomsection\label{I-elemprop}

\begin{enumerate}
\def\labelenumi{\arabic{enumi}.}
\tightlist
\item
  \(\forall\, A \in \mathcal{A}\), \(\mathbb{P}(A) \in [0,1]\) et
  \(\mathbb{P}(A^c)= 1-\mathbb{P}(A)\).
\item
  \(\forall\, A,B \in \mathcal{A}\),
  \(A \subset B \Rightarrow \mathbb{P}(A) \leq \mathbb{P}(B)\).
\item
  \(\forall\, A,B \in \mathcal{A}\),
  \(\mathbb{P}(A \cup B ) = \mathbb{P}(A) + \mathbb{P}(B) - \mathbb{P}(A \cap B)\).
\item
  Inégalité de Boole~: \(\forall\, n \in \mathbb{N}^\ast\),
  \(\forall\, (A_i)_{1 \leq i \leq n} \in \mathcal{A}\),
  \[\mathbb{P}\left(\bigcup_{i=1}^n A_i\right) \leq \sum_{i=1}^n \mathbb{P}(A_i).\]
\item
  Formule de Poincaré~: \(\forall\, n \in \mathbb{N}^\ast\),
  \(\forall\, (A_i)_{1 \leq i \leq n} \in \mathcal{A}\)
  \[ \mathbb{P}\left(\bigcup_{i=1}^n A_i\right) = \sum_{i=1}^n \mathbb{P}(A_i) - \sum_{1 \leq i < j \leq n} \mathbb{P}(A_i \cap A_j) + \ldots + (-1)^n \mathbb{P}\left(\bigcap_{i=1}^n A_i\right).\]
\end{enumerate}
\end{proposition}

\runinhead{Exercice -- Démonstration (\(\mathord{\bullet}\))}\phantomsection\addcontentsline{toc}{subsubsection}{Démonstration}\label{I-propelem}
\solref{I-answer-propelem}

\addcontentsline{toc}{subsection}{Théorème de la continuité monotone}
\begin{theoreme}[{Théorème de la continuité
monotone}]\phantomsection\label{I-continuitemonotone}

Dans le cas d'une suite \((A_n)_{n\in\mathbb{N}^\ast}\) d'éléments de
\(\mathcal{A}\) croissante, on a
\[ \mathbb{P}\left(\bigcup_{n\in\mathbb{N}^\ast} A_n\right) = \lim_{n \rightarrow \infty} \mathbb{P}(A_n).\]
\end{theoreme}

\runinhead{Exercice -- Démonstration (\(\mathord{\bullet}\))}\phantomsection\addcontentsline{toc}{subsubsection}{Démonstration}\label{I-contmon}
\solref{I-answer-contmon}

\runinhead{Exercice -- Une définition alternative de la
probabilité (\(\mathord{\bullet}\))}\phantomsection\addcontentsline{toc}{subsubsection}{Une
définition alternative de la probabilité}\label{I-altdef}
Soit \((\Omega, \mathcal{A})\) un espace probabilisable. Supposons que
\(\mathbb{P}: \mathcal{A}\to [0,1]\) vérifie~:

\begin{enumerate}
\def\labelenumi{\arabic{enumi}.}
\tightlist
\item
  \(\mathbb{P}(\Omega) = 1\),
\item
  Pour \(A, B \in \mathcal{A}\), tels que \(A\cap B = \varnothing\),
  \(\mathbb{P}(A\cup B)= \mathbb{P}(A) + \mathbb{P}(B)\) (additivité),
\item
  Pour toute suite \((A_n)_{n\in\mathbb{N}^\ast}\) d'éléments de
  \(\mathcal{A}\) croissante
  \[\mathbb{P}\left(\bigcup_{n\in\mathbb{N}^\ast} A_n\right) = \lim_{n \rightarrow \infty} \mathbb{P}(A_n).\]
\end{enumerate}

Montrer que \(\mathbb{P}\) vérifie la propriété de
\protect\hyperref[I-defproba]{\(\sigma\)-additivité (p.~\pageref*{I-defproba})}.\solref{I-answer-altdef}

\addcontentsline{toc}{subsection}{Suite décroissante}
\begin{remarque}[{Suite
décroissante}]\phantomsection\label{I-suite-duxe9croissante}

Dans le cas d'une suite décroissante, on a
\[ \mathbb{P}\left(\bigcap_{n \in \mathbb{N}^\ast} A_n\right) = \lim_{n \rightarrow \infty} \mathbb{P}(A_n).\]

Le second point de la \protect\hyperref[I-defproba]{définition de la
probabilité (p.~\pageref*{I-defproba})} donne la probabilité de la
réunion \(\cup_n A_n\) en fonction des \(\mathbb{P}(A_n)\) lorsque les
événements sont deux à deux disjoints. Si ce n'est pas le cas, on a tout
de même la majoration suivante~:
\end{remarque}

\addcontentsline{toc}{subsection}{Une majoration bien utile}
\begin{proposition}[{Une majoration bien
utile}]\phantomsection\label{I-une-majoration-bien-utile}

Soit \(\mathbb{P}\) une probabilité et soit
\((A_n)_{n\in \mathbb{N}^\ast}\) une famille dénombrable d'événements.
On a alors
\[ \mathbb{P}(\cup_n A_n) \leq \sum_{n\in\mathbb{N}^\ast} \mathbb{P}(A_n) \]
\end{proposition}

\runinhead{Exercice -- Démonstration (\(\mathord{\bullet}\))}\phantomsection\addcontentsline{toc}{subsubsection}{Démonstration}\label{I-boolinf}
\solref{I-answer-boolinf}

\section{Probabilité
conditionnelle}\label{I-probabilituxe9-conditionnelle}

La construction d'un modèle probabiliste repose sur l'information connue
\textbf{a priori} sur l'expérience aléatoire. Ce modèle permet de
quantifier les probabilités de réalisation de certains résultats de
l'expérience. Il est fondamental de remarquer que si l'information
change, les probabilités de réalisation changent.

\addcontentsline{toc}{subsection}{Information a priori}
\begin{exemple}[{Information a
priori}]\phantomsection\label{I-information-a-priori}

On cherche pour un lancer de deux dés, la probabilité de l'événement
``la somme est supérieure ou égale à 10''. Elle vaut 1/6 sans
information supplémentaire, 1/2 si l'on sait que le résultat d'un des
dés est 6, 0 si l'on sait a priori que le résultat d'un des dés est 2.
Pour obtenir ces résultats, on a calculé dans chaque cas le rapport du
nombre de résultats favorables sur le nombre de cas possibles. Il est
ainsi indispensable de bien définir l'espace de probabilité lié à
l'expérience munie de l'information a priori. On remarque également que
l'information a priori a changé la valeur de la probabilité de
l'événement.

L'outil qui va nous permettre d'introduire de l'information est la
probabilité conditionnelle dont nous donnons ici la définition.
\end{exemple}

\addcontentsline{toc}{subsection}{Probabilité conditionnelle}
\begin{definition}[{Probabilité
conditionnelle}]\phantomsection\label{I-defprobacond}

Soient \((\Omega, \mathcal{A}, \mathbb{P})\) un espace probabilisé,
\(A, B \in \mathcal{A}\) tels que \(\mathbb{P}(B)>0\). La
\emph{probabilité conditionnelle} de \(A\) sachant \(B\) est le nombre
\begin{equation*}
\mathbb{P}(A|B) = \frac{\mathbb{P}(A\cap B)}{\mathbb{P}(B)}.
\end{equation*}
\end{definition}

Cela définit une probabilité comme le montre la proposition suivante.

\addcontentsline{toc}{subsection}{Conséquences}
\begin{proposition}[{Conséquences}]\phantomsection\label{I-consuxe9quences}

\begin{enumerate}
\def\labelenumi{\arabic{enumi}.}
\tightlist
\item
  Soient \((\Omega, \mathcal{A}, \mathbb{P})\) un espace probabilisé et
  \(B \in \mathcal{A}\) tel que \(\mathbb{P}(B)>0\). Alors l'application
  de \(\mathcal{A}\) dans \([0, 1]\) qui à \(A\) associe
  \(\mathbb{P}(A|B)\) définit une nouvelle probabilité sur \(\Omega\),
  appelée probabilité conditionnelle sachant \(B\).
\item
  Si \(\mathbb{P}(A) > 0\) et \(\mathbb{P}(B) > 0\) , nous avons
  \[\mathbb{P}(A|B)\, \mathbb{P}(B) = \mathbb{P}(A \cap B) = \mathbb{P}(B|A)\, \mathbb{P}(A).\]
\end{enumerate}
\end{proposition}

\begin{preuve}
Il est clair que \(0 \leq \mathbb{P}(A|B) \leq 1\). Par ailleurs, les
deux propriétés de la \protect\hyperref[I-defproba]{définition de la
probabilité (p.~\pageref*{I-defproba})} pour \(\mathbb{P}(\cdot|B)\)
proviennent des mêmes propriétés pour \(\mathbb{P}\) et des remarques
suivantes~: \(\Omega \cap B = B\), et
\((\bigcup_{n\in\mathbb{N}^\ast} A_n ) \cap B = \bigcup_{n\in\mathbb{N}^\ast} (A_n \cap B)\).
De plus, si \(A\) et \(C\) sont disjoints, il en est de même de
\(A \cap B\) et \(C \cap B\). L'assertion 2 est évidente, d'après la
définition de la \protect\hyperref[I-defprobacond]{Probabilité
conditionnelle (p.~\pageref*{I-defprobacond})}.\hfill$\blacksquare$
\end{preuve}

\addcontentsline{toc}{subsection}{Formule des probabilités totales}
\begin{proposition}[{Formule des probabilités
totales}]\phantomsection\label{I-formprobatot}

Soit \((B_n)_{n\in\mathbb{N}^\ast}\) une partition finie ou dénombrable
d'événements de \(\Omega\) (i.e.~telle que
\(\bigcup_{n\in\mathbb{N}^\ast} B_n = \Omega\) et les \(B_n\) sont
deux-à-deux disjoints), telle que \(\mathbb{P}(B_n ) > 0\) pour tout
\(n\in\mathbb{N}^\ast\). Pour tout \(A \in \mathcal{A}\), on a alors
\begin{equation*}
P(A) = \sum_{n\in\mathbb{N}^\ast} P(A \cap B_n ) = \sum_{n\in\mathbb{N}^\ast} P(A|B_n)\, P(B_n).
\end{equation*}
\end{proposition}

\begin{preuve}
Nous avons \(A = \bigcup_{n\in\mathbb{N}^\ast} (A\cap B_n)\). Par
hypothèse, les ensembles \((A\cap B_n)\) sont deux-à-deux disjoints et
de plus \(\mathbb{P}(A\cap B_n) = \mathbb{P}(A|B_n)\,\mathbb{P}(B_n)\).
Le résultat découle du deuxième point de la
\protect\hyperref[I-defproba]{définition de la
probabilité (p.~\pageref*{I-defproba})}.\hfill$\blacksquare$
\end{preuve}

\addcontentsline{toc}{subsection}{Formule de Bayes}
\begin{proposition}[{Formule de Bayes}]\phantomsection\label{I-bayes}

Selon les mêmes hypothèses que ci-dessus et si \(\mathbb{P}(A) > 0\), on
a

\begin{equation*}
\forall\, i \in \{1,\dots,n\},\  \mathbb{P}(B_i | A) = \dfrac{\mathbb{P}(A | B_i)\, \mathbb{P}(B_i)}{\sum_{n\in\mathbb{N}^\ast} \mathbb{P}(A | B_n)\, \mathbb{P}(B_n)}.
\end{equation*}
\end{proposition}

\begin{preuve}
Le dénominateur vaut \(\mathbb{P}(A)\) d'après la
\protect\hyperref[I-formprobatot]{Formule des probabilités
totales (p.~\pageref*{I-formprobatot})}. La définition de la
\protect\hyperref[I-defprobacond]{probabilité
conditionnelle (p.~\pageref*{I-defprobacond})} implique~:
\begin{equation*}
\mathbb{P}(B_i | A) = \dfrac{\mathbb{P}(A \cap B_i)}{\mathbb{P}(A)} = \dfrac{\mathbb{P}(A | B_i) \mathbb{P}(B_i)}{\mathbb{P}(A)}.
\end{equation*}\hfill$\blacksquare$
\end{preuve}

\begin{remarque}\phantomsection\label{I-remarque}
La formule de Bayes, simple conséquence des axiomes et de la définition
de la probabilité conditionnelle, tient une place à part dans le calcul
des probabilités en raison de son importance pratique considérable et
des controverses auxquelles son application pratique a donné lieu~: elle
est à la base de toute une branche de la statistique appelée statistique
bayésienne.
\end{remarque}

\runinhead{Exercice -- Test (\(\mathord{\bullet}\))}\phantomsection\addcontentsline{toc}{subsubsection}{Test}\label{I-test}
Un individu est tiré au hasard dans une population où l'on trouve une
proportion \(10^{-4}\) de séropositifs. On lui fait passer un test de
détection de la séropositivité. Par ailleurs, des expérimentations
antérieures ont permis de savoir que les probabilités d'avoir un
résultat positif lors de l'application du test si l'individu est
séropositif, ou s'il ne l'est pas, sont respectivement égales à 0,99
(c'est la sensibilité du test) et à 0,001 (0,999 = 1 - 0,001 est la
spécificité du test). Sachant que le test donne un résultat positif,
quelle est la probabilité pour que l'individu soit effectivement
séropositif\,?\solref{I-answer-test}

\section{Indépendance des
événements}\label{I-induxe9pendance-des-uxe9vuxe9nements}

La notion d'indépendance est absolument fondamentale en probabilités et
nous verrons par la suite toutes ses implications dans la modélisation
de l'aléatoire.

Intuitivement, deux événements \(A\) et \(B\) sont indépendants si le
fait de savoir que \(A\) est réalisé ne donne aucune information sur la
réalisation de \(B\) et réciproquement.

Si \(B\) est un événement de probabilité strictement positive, \(A\)
sera dit indépendant de \(B\) si
\[\mathbb{P}(A | B) = \frac{\mathbb{P}(A\cap B)}{\mathbb{P}(B)} = \mathbb{P}(A)\]
On remarque que cette formule se symétrise et la notion d'indépendance
se définit finalement comme suit.

\addcontentsline{toc}{subsection}{Indépendance de deux événements}
\begin{definition}[{Indépendance de deux
événements}]\phantomsection\label{I-induxe9pendance-de-deux-uxe9vuxe9nements}

Deux événements \(A\) et \(B\) sont \emph{indépendants} si et seulement
si \begin{equation*}
\mathbb{P}(A\cap B) = \mathbb{P}(A)\, \mathbb{P}(B).
\end{equation*}
\end{definition}

\begin{remarque}\phantomsection\label{I-remarques}
\begin{itemize}
\tightlist
\item
  La probabilité de voir \(A\) réalisé ne dépend pas de la réalisation
  de \(B\), et réciproquement.
\item
  Cette notion est liée au choix de la probabilité \(\mathbb{P}\) et
  n'est pas une notion ensembliste. Cela n'a en particulier rien à voir
  avec le fait que \(A\) et \(B\) soient disjoints ou non.
\item
  Si \(\mathbb{P}(A)>0\) et \(\mathbb{P}(B)>0\), alors \begin{equation*}
   \mathbb{P}(A\cap B) = \mathbb{P}(A)\, \mathbb{P}(B) \Leftrightarrow \mathbb{P}(A) = \mathbb{P}(A|B) \Leftrightarrow \mathbb{P}(B) = \mathbb{P}(B|A).
   \end{equation*}
\end{itemize}
\end{remarque}

\begin{proposition}\phantomsection\label{I-proposition}
Si les événements \(A\) et \(B\) sont indépendants, alors il en est de
même des couples \((A^c,B)\), \((A,B^c)\) et \((A^c,B^c)\).
\end{proposition}

\runinhead{Exercice -- Démonstration (\(\mathord{\bullet}\))}\phantomsection\addcontentsline{toc}{subsubsection}{Démonstration}\label{I-indep}
\solref{I-answer-indep}

\runinhead{Exercice -- Auto-indépendant\,? (\(\mathord{\bullet}\))}\phantomsection\addcontentsline{toc}{subsubsection}{Auto-indépendant\,?}\label{I-autoindep}
À quelle condition un événement \(A\) est-il indépendant de
lui-même\,?\solref{I-answer-autoindep}

\begin{exemple}\phantomsection\label{I-exemples-4}
\begin{enumerate}
\def\labelenumi{\arabic{enumi}.}
\tightlist
\item
  On lance 3 fois un dé. Si \(A_i\) est un événement qui ne dépend que
  du \(i^\text{ème}\) lancer, alors \(A_1\) , \(A_2\) , \(A_3\) sont
  indépendants.
\item
  On tire une carte au hasard dans un jeu de 52 cartes. Soit \(A\) =
  \{la carte est une dame\} et \(B\) = \{la carte est un coeur\}. Il est
  facile de voir que \(\mathbb{P}(A) = 4/52\) et
  \(\mathbb{P}(B) = 13/52\) et
  \(\mathbb{P}(A \cap B) = \mathbb{P}( \{\text{la carte est la dame de coeur}\} ) = 1/52 = \mathbb{P}(A) \mathbb{P}(B)\).
  Ainsi, les événements \(A\) et \(B\) sont indépendants pour la
  probabilité uniforme \(\mathbb{P}\).
\item
  On suppose maintenant que le jeu de cartes soit trafiqué. Soit
  \(\widetilde{\mathbb{P}}\) la nouvelle probabilité correspondant au
  tirage de cartes. On suppose également que
  \[ \widetilde{\mathbb{P}}(\{\text{As de trèfle}\} ) = \frac{1}{2}, \,\,\,\,\, \widetilde{\mathbb{P}}(\{\text{autre carte}\}) = \frac{1}{2} \frac{1}{51} = \frac{1}{102},\]
  alors
  \[\widetilde{\mathbb{P}}(A \cap B) = \frac{1}{102} \neq \widetilde{\mathbb{P}}(A)\widetilde{\mathbb{P}}(B) = \frac{2}{51}\frac{13}{102}.\]
  Les événements \(A\) et \(B\) ne sont pas indépendants sous la
  probabilité \(\widetilde{\mathbb{P}}\).
\end{enumerate}
\end{exemple}

\addcontentsline{toc}{subsection}{Indépendance de \(n\) événements}
\begin{definition}[{Indépendance de \(n\)
événements}]\phantomsection\label{I-induxe9pendance-de-n-uxe9vuxe9nements}

\(n\) événements \(A_1, \dots, A_n\) sont \emph{mutuellement
indépendants} ou simplement \emph{indépendants} si et seulement si pour
toute partie \(I \subset \lbrace 1,\dots, n\rbrace\) \begin{equation*}
\mathbb{P}(\cap_{i\in I} A_i) = \prod_{i\in I} \mathbb{P}(A_i).
\end{equation*}
\end{definition}

\begin{remarque}\phantomsection\label{I-remarque-1}
Attention, l'indépendance des événements deux à deux n'implique pas
l'indépendance mutuelle des événements.
\end{remarque}

\section{Réflexion sur le concept de
probabilité}\label{I-remarque-ruxe9flexion-sur-le-concept-de-probabilituxe9}

La théorie mathématiques des probabilités ne dit pas quelle loi de
probabilité choisir sur un espace \((\Omega,\mathcal{A})\) parmi
l'infinité de lois possibles. Ce problème qui concerne ceux qui veulent
appliquer le calcul des probabilités, renvoie à la nature ``physique''
du concept de probabilité qui formalise et quantifie le sentiment
d'incertitude vis-à-vis d'un événement. Ce problème d'ordre conceptuel
oppose deux écoles de pensée, la conception objectiviste et la
conception subjectiviste.

Pour les tenants du premier point de vue, la probabilité d'un événement
peut être déterminée de manière unique. Dans la vision dite classique,
héritée des jeux de hasard, \(\Omega\) est fini et on donne à chaque
événement élémentaire la même probabilité. Le calcul des probabilités se
résume alors à un problème de dénombrement et la probabilité d'un
événement est le rapport du nombre de cas favorables sur le nombre de
cas possibles. Dans le cas infini, la vision fréquentiste repose sur la
loi des grands nombres~: si on répète un grand nombre de fois
l'expérience, la proportion de fois où un événement sera réalisé va
converger vers la probabilité de cet événement. Dans ce cadre, il est
impossible de donner une valeur et même un sens à un événement non
répétable comme ``pleuvra-t-il demain\,?''. En outre, la répétition à
l'infini d'une même expérience étant physiquement irréalisable, la loi
des grands nombres étant un résultat qui suppose défini le concept de
probabilité, la vision fréquentiste est logiquement intenable.

Dans la conception subjectiviste, la probabilité objective d'un
événement n'existe pas et n'est donc pas une grandeur mesurable analogue
à la masse d'un corps, par exemple. C'est simplement une mesure
d'incertitude qui reflète un degré de croyance pouvant varier avec les
circonstances et l'observateur, donc subjective, la seule exigence étant
qu'elle satisfasse aux axiomes du calcul des probabilités. Des méthodes
ont alors été proposées pour passer d'un simple pré-ordre sur les
événements, à une probabilité. Puisque la répétition n'est plus
nécessaire, on peut probabiliser des événements non répétables et
étendre ainsi le domaine d'application du calcul des probabilités,
notamment pour orienter des prises de décisions. On notera que la
\protect\hyperref[I-bayes]{formule de Bayes (p.~\pageref*{I-bayes})}
permet d'intégrer facilement de l'information a priori, dans la mesure
où celle-ci est probabilisée.

On arrête ici ces quelques remarques sans prendre parti dans une
querelle qui dure encore. L'un ou l'autre point de vue sera adopté selon
les ouvrages rencontrés. Dans tous les cas, les outils mathématiques
développés dans ce cours seront adaptés. On rappelle tout de même que la
modélisation probabiliste a prouvé son efficacité dans de nombreuses
applications mais que, comme tout modèle, ce n'est qu'une représentation
simplificatrice de la réalité et que ses hypothèses doivent être mises à
l'épreuve des faits. À ce titre, on citera Georges Matheron qui dans son
essai sur la pratique des probabilités Estimer et Choisir (Matheron
(\protect\hyperlink{refI-matheron}{1978})) écrit fort justement~: ``Il
n'y a pas de probabilités en soi. Il n'y a que des modèles
probabilistes''.

\chapter{Variables aléatoires}\label{I-variables-aluxe9atoires-1}

En théorie moderne des probabilités, on préfère prendre un point de vue
fonctionnel plutôt qu'ensembliste, et utiliser les variables aléatoires
plutôt que les événements. Une variable aléatoire est une grandeur qui
dépend du résultat de l'expérience. Par exemple,

\begin{itemize}
\tightlist
\item
  le nombre de 6 obtenus dans un lancer de 3 dés,
\item
  le nombre d'appels dans un central téléphonique pendant une heure,
\item
  la distance du point d'atteinte d'une flèche au centre de la cible,
\item
  la valeur maximale d'un prix d'actif sur un intervalle de temps donné,
\end{itemize}

sont des variables aléatoires.

\addcontentsline{toc}{subsection}{Variable aléatoire}
\begin{definition}[{Variable
aléatoire}]\phantomsection\label{I-variable-aluxe9atoire}

Soient \((\Omega,\mathcal{A})\) et \((E,\mathcal{E})\) deux espaces
probabilisables (mesurables). Une \emph{variable aléatoire} \(X\) à
valeurs dans \(E\) est une application mesurable de
\((\Omega,\mathcal{A})\) dans \((E,\mathcal{E})\), c'est-à-dire telle
que \(\forall B \in \mathcal{E}\), \(X^{-1}(B) \in \mathcal{A}\).
\end{definition}

En pratique, l'ensemble \(E\) pourra être un ensemble fini ou
dénombrable ou \(\mathbb{R}\) ou \(\mathbb{R}^d\) ou encore un espace
plus sophistiqué tel que l'ensemble \(C(\mathbb{R}_+ , \mathbb{R}^d)\)
des fonctions continues de \(\mathbb{R}_+\) dans \(\mathbb{R}^d\).

\addcontentsline{toc}{subsection}{Vocabulaire}
\begin{remarque}[{Vocabulaire}]\phantomsection\label{I-vocabulaire}

La terminologie, consacrée par l'usage, peut être trompeuse. Une
variable aléatoire n'est pas une variable (au sens de l'analyse) mais
une fonction. Cette terminologie est apparentée à la notion de variable
en physique ou en sciences humaines où on désigne volontiers par
``variable'' la valeur prise par une fonction de l'état du système
étudié.

L'intérêt principal de travailler avec des variables aléatoires est de
pouvoir substituer à l'espace abstrait \(\Omega\) des résultats de
l'expérience l'espace \(E\), mieux connu dans la pratique. Ainsi, grâce
à une variable aléatoire \(X\), nous pouvons transporter la structure
abstraite du modèle probabiliste \((\Omega, \mathcal{A}, \mathbb{P})\)
sur l'espace d'arrivée \(E\), en posant pour \(B \in \mathcal{E}\)
\[\mathbb{P}_X (B) = \mathbb{P}(X^{-1}(B)) = \mathbb{P}(\{\omega, X(\omega)\in B\})\]

Cette formule définit une nouvelle probabilité, notée \(\mathbb{P}_X\)
et définie sur \((E,\mathcal{E})\), qui s'appelle la \emph{loi de la
variable} \(X\). On dit aussi que \(\mathbb{P}_X\) est la \emph{mesure
(de probabilité) image} de \(\mathbb{P}\) par \(X\).

Pour se convaincre que \(\mathbb{P}_X\) définit bien une probabilité sur
\((E,\mathcal{E})\), il suffit de remarquer que \begin{align*}
& X^{-1}(\varnothing) = \varnothing, X^{-1}(E) = \Omega, X^{-1}(B^c) = X^{-1}(B)^c \\
& X^{-1}(\cap_i A_i) = \cap_i X^{-1}(A_i), X^{-1}(\cup_i A_i) = \cup_i X^{-1}(A_i)
\end{align*}
\end{remarque}

\addcontentsline{toc}{subsection}{Notation}
\begin{remarque}[{Notation}]\phantomsection\label{I-notation}

Il est usuel de noter l'ensemble
\(X^{-1}(B) = \{\omega \in \Omega, X(\omega) \in B\}\) par
\(\{X \in B\}\), ce qui allège les écritures. On se rappelera néanmoins
que cette notation désigne un sous-ensemble de \(\Omega\).
\end{remarque}

\(\mathbb{P}_X\) sera plus facile à caractériser que \(\mathbb{P}\)
puisque \(E\) est un ensemble connu (on pourra en particulier utiliser
ses propriétés topologiques) alors que \(\Omega\) est un espace
abstrait. Les variables que nous rencontrerons dans ce cours seront soit
à valeurs dans un ensemble dénombrable, soit à valeurs dans
\(\mathbb{R}\) ou dans \(\mathbb{R}^d\). Nous les appellerons
respectivement des variables aléatoires discrètes, réelles ou des
vecteurs aléatoires. Leurs lois seront alors des probabilités
respectivement sur un ensemble dénombrable, sur \(\mathbb{R}\) ou sur
\(\mathbb{R}^d\). Le cas discret est considéré connu.

\section{Variables aléatoires
réelles}\label{I-variables-aluxe9atoires-ruxe9elles}

Dans le cas où \(E = \mathbb{R}\), on travaille généralement avec la
\emph{tribu des boréliens}, que l'on notera \(\mathcal{B}(\mathbb{R})\),
qui est très utile comme on le verra pour identifier les lois de
probabilité. On rappelle ici quelques propriétés sur les tribus, déjà
vues dans le cours d'intégration.

\addcontentsline{toc}{subsection}{Tribu engendrée}
\begin{definition}[{Tribu
engendrée}]\phantomsection\label{I-tribu-engendruxe9e}

Si \(C \subset \mathcal{P}(\Omega)\), on appelle tribu engendrée par
\(C\) la plus petite tribu contenant \(C\). Elle existe toujours, car
d'une part \(\mathcal{P}(\Omega)\) est une tribu contenant \(C\), et
d'autre part l'intersection d'une famille quelconque de tribus est une
tribu. Ainsi, la tribu engendrée par \(C\) est l'intersection de toutes
les tribus contenant \(C\).
\end{definition}

\addcontentsline{toc}{subsection}{Tribus engendrées}
\begin{exemple}[{Tribus
engendrées}]\phantomsection\label{I-tribus-engendruxe9es}

\begin{itemize}
\tightlist
\item
  La tribu engendrée par un ensemble \(A \subset \Omega\) est
  \({\varnothing, A, A^c , \Omega}\).
\item
  si \((A_i )_{i \in I}\) est une partition finie ou dénombrable de
  \(\Omega\) (i.e.~les \(A_i\) sont deux-à-deux disjoints et leur
  réunion est \(\Omega\)), la tribu engendrée par \(\{A_i , i \in I\}\)
  est l'ensemble des réunions \(B_J = \cup_{i \in J} A_i\), où \(J\)
  décrit la classe de toutes les parties de I.
\end{itemize}
\end{exemple}

\addcontentsline{toc}{subsection}{Tribu borélienne sur \(\mathbb{R}\)}
\begin{definition}[{Tribu borélienne sur
\(\mathbb{R}\)}]\phantomsection\label{I-bor}

Si \(\Omega = \mathbb{R}\), on appelle tribu borélienne, que l'on note
\(\mathcal{B}(\mathbb{R})\), la tribu engendrée par la classe des
ouverts de \(\mathbb{R}\).
\end{definition}

À titre d'exercice de maniement des tribus, on donne en détail la
démonstration du résultat suivant~:

\addcontentsline{toc}{subsection}{Définition alternative}
\begin{proposition}[{Définition
alternative}]\phantomsection\label{I-altbor}

La tribu borélienne de \(\mathbb{R}\) est la tribu engendrée par les
intervalles de la forme \(] - \infty, a]\) pour \(a \in \mathbb{Q}\).
\end{proposition}

\begin{preuve}
On rappelle que toute tribu est stable par passage au complémentaire,
par réunion ou intersection dénombrable.

Soit \(C_1\) la classe des intervalles ouverts de \(\mathbb{R}\),
\(C_2\) la classe des intervalles \(] - \infty, a]\) pour
\(a \in \mathbb{Q}\) et \(\mathcal{B}(\mathbb{R})\) la tribu borélienne.

Tout ouvert \(A\) est réunion dénombrable d'intervalles ouverts, en
effet, on a \(A = \cup_{(q,n)\in B} ]q -\frac{1}{n}, q -\frac{1}{n}[\),
où \(B\) est l'ensemble (dénombrable) des couples \((q,n)\) avec
\(q \in \mathbb{Q}\), \(n \in \mathbb{N}\) et
\(]q -\frac{1}{n}, q -\frac{1}{n}[ \in A\). La tribu engendrée par les
intervalles ouverts est donc identique à la tribu borélienne par la
stabilité par réunion dénombrable et le fait que
\(C_1 \subset \mathcal{B}(\mathbb{R})\).

Réciproquement, soit \(]x, y[\) un intervalle ouvert de \(\mathbb{R}\).
Soit \((x_n)_n\) une suite de rationnels décroissant vers \(x\) et
\((y_n)_n\) une suite de rationnels croissant strictement vers y. On a~:
\[ ]x, y[= \bigcup_n (] - \infty, y_n ]\bigcap ] - \infty, x_n ]^c ).\]
Donc \(C_1\) et par conséquent \(\mathcal{B}(\mathbb{R})\) sont inclus
dans la tribu engendrée par \(C_2\). Finalement, comme tout fermé est
dans \(\mathcal{B}(\mathbb{R})\) par complémentarité, on a aussi
\(C_2 \subset \mathcal{B}(\mathbb{R})\), d'où le
résultat.\hfill$\blacksquare$
\end{preuve}

On peut maintenant définir une variable aléatoire réelle~:

\addcontentsline{toc}{subsection}{Variable aléatoire réelle}
\begin{definition}[{Variable aléatoire
réelle}]\phantomsection\label{I-defvar}

Soit l'espace d'état \(\Omega\) muni de la tribu \(\mathcal{A}\) des
événements. Une application \(X\) de \(\Omega\) dans \(\mathbb{R}\) est
une \emph{variable aléatoire réelle} si \(X^{-1}(B) \in \mathcal{A}\)
pour tout \(B \in \mathcal{B}(\mathbb{R})\), autrement dit si \(X\) est
une application mesurable de \((\Omega, \mathcal{A})\) dans
\((\mathbb{R},\mathcal{B}(\mathbb{R}))\).
\end{definition}

\addcontentsline{toc}{subsection}{Loi d'une variable aléatoire réelle}
\begin{definition}[{Loi d'une variable aléatoire
réelle}]\phantomsection\label{I-defloivar}

La probabilité \(\mathbb{P}_X\), définie sur
\((\mathbb{R},\mathcal{B}(\mathbb{R}))\) par
\(\mathbb{P}_X (B) = \mathbb{P}(X^{-1}(B))\) pour
\(B \in \mathcal{B}(\mathbb{R})\) est appelée \emph{loi de la variable
\(X\)}, ou \emph{distribution} de \(X\).
\end{definition}

On peut voir \(\mathbb{P}_X\) comme une transposition de \(\mathbb{P}\)
sur \(\mathbb{R}\). On a alors le résultat très utile suivant~:

\addcontentsline{toc}{subsection}{Composition}
\begin{proposition}[{Composition}]\phantomsection\label{I-composition}

Si \(X_1, \ldots, X_n\) sont des variables aléatoires réelles et si
\(g\) est une fonction mesurable de \(\mathbb{R}^n\) dans \(\mathbb{R}\)
(munis de leur tribu borélienne), alors \(Y = g(X_1,\ldots,X_n)\) est
une variable aléatoire réelle.
\end{proposition}

\begin{preuve}
On note \(X = (X_1, \ldots, X_n)\). \(X\) est mesurable de
\((\Omega,\mathcal{A})\) dans
\((\mathbb{R}^n,\mathcal{B}(\mathbb{R}^n))\) puisque ce vecteur est
constitué de composantes mesurables de \((\Omega,\mathcal{A})\) dans
\((\mathbb{R},\mathcal{B}(\mathbb{R}))\) et que
\(\mathcal{B}(\mathbb{R}^n) = \mathcal{B}(\mathbb{R})\otimes \ldots \otimes \mathcal{B}(\mathbb{R})\).
Soit \(B \in \mathcal{B}(\mathbb{R})\), alors
\(Y^{-1}(B) = (g \circ X)^{-1}(B) = X^{-1}(g^{-1}(B))\). Puisque \(g\)
et \(X\) sont mesurables, \(g^{-1}(B) \in \mathcal{B}(\mathbb{R}^n)\) et
\(X^{-1}(g^{-1}(B))\in \mathcal{A}\).\hfill$\blacksquare$
\end{preuve}

Comme application de ce résultat, on a les propriétés suivantes~:

\addcontentsline{toc}{subsection}{Conséquences}
\begin{proposition}[{Conséquences}]\phantomsection\label{I-consuxe9quences-1}

Soient \(X\), \(Y\) et \((X_n)_{n \in \mathbb{N}^\star}\) des variables
aléatoires réelles. On a

\begin{enumerate}
\def\labelenumi{\arabic{enumi}.}
\item
  \(X + Y\), \(XY\), \(\frac{X}{Y}\) si \(Y \neq 0\), sont des variables
  aléatoires.
\item
  \(\sup_{1\leq p \leq n} X_p\), \(\inf_{1\leq p \leq n} X_p\), sont des
  variables aléatoires.
\item
  \(\sup_{n\geq 1} X_n\), \(\inf_{n\geq 1} X_n\), sont des variables
  aléatoires.
\item
  Si \(X_n(\omega) \xrightarrow[n \to \infty]{} Z(\omega)\),
  \(\forall \omega\), alors la limite \(Z\) est une variable aléatoire.
\item
  \(Z = 1_A\) est une variable aléatoire \(\Leftrightarrow\)
  \(A \in \mathcal{A}\).
\end{enumerate}
\end{proposition}

\section{Loi des variables aléatoires
réelles}\label{I-loi-des-variables-aluxe9atoires-ruxe9elles}

Nous avons vu précédemment la définition générale d'une probabilité
\(\mathbb{P}\) sur un espace quelconque \(\Omega\) muni d'une tribu
\(\mathcal{A}\). Un problème fondamental est de construire et de
caractériser ces probabilités. La résolution de ce problème lorsque
\(\Omega\) est fini ou dénombrable est supposée connue. Puisque les
variables aléatoires réelles sont définies sur \(\mathbb{R}\), nous
allons étudier leur loi de probabilité \(\mathbb{P}_X\) définie sur
\(\mathbb{R}\) muni de la tribu \(\mathcal{B}(\mathbb{R})\).

\addcontentsline{toc}{subsection}{Fonction de répartition}
\begin{definition}[{Fonction de
répartition}]\phantomsection\label{I-deffdr}

Soit \(X\) une variable aléatoire réelle et \(\mathbb{P}_X\) sa loi. La
\emph{fonction de répartition} de \(X\) est la fonction
\begin{equation*}
F_X(x) = \mathbb{P}_X(\left]-\infty, x\right]) = \mathbb{P}(X \leq x),\ x \in \mathbb{R}.
\end{equation*}
\end{definition}

\addcontentsline{toc}{subsection}{La fonction de répartition caractérise
la probabilité}
\begin{theoreme}[{La fonction de répartition caractérise la
probabilité}]\phantomsection\label{I-carac}

La fonction de répartition \(F\) caractérise la probabilité
\(\mathbb{P}_X\) sur (\(\mathbb{R},\mathcal{B}(\mathbb{R})\)).
\end{theoreme}

\begin{preuve}
On veut montrer que la connaissance de \(F\) telle que définie
\protect\hyperref[I-deffdr]{ci-dessus (p.~\pageref*{I-deffdr})}
détermine \(\mathbb{P}_X\) de manière unique, c'est-à-dire que si il
existe une autre probabilité \(\mathbb{P}_Y\) telle que
\(G(x) = \mathbb{P}_Y(]-\infty,x])\) pour \(x \in \mathbb{R}\) et si
\(F = G\), alors \(\mathbb{P}_X = \mathbb{P}_Y\).

Soit \(\mathcal{B}_0\) l'ensemble de toutes les réunions finies
d'intervalles disjoints de la forme \(]x,y]\) ou
\(-\infty \leq x < y \leq +\infty\). Cet ensemble contient
\(\mathbb{R}\), \(\varnothing\) et est stable par passage au
complémentaire et par réunion finie, c'est donc une algèbre. De manière
évidente, il est aussi stable par intersection finie et est ainsi un
\(\pi\)-système. Par ailleurs, on a déjà vu que
\(\sigma(\mathcal{B}_0)\) (la tribu engendrée par \(\mathcal{B}_0\)) est
la tribu borélienne \(\mathcal{B}(\mathbb{R})\).

D'après \protect\hyperref[I-deffdr]{la définition de la fonction de
répartition (p.~\pageref*{I-deffdr})}, on a
\(\mathbb{P}_X(]x,y]) = F(y)-F(x)\) pour tous \(x< y\). Par conséquent,
si \(B = \cup_{i=1}^n ]x_i,y_i]\), avec \(x_i < y_i < x_{i+1}\), on a
\[\mathbb{P}_X(B) = \sum_{i=1}^n \mathbb{P}_X(]x_i,y_i]) = \sum_{i=1}^n F(y_i)-F(x_i),\]
car les intervalles sont disjoints.

Si \(\mathbb{P}_Y\) est une autre probabilité telle que
\(F(x) = \mathbb{P}_Y(]-\infty,x])\), alors ce qui précède montre que
\(\mathbb{P}_X = \mathbb{P}_Y\) sur \(\mathcal{B}_0\) qui engendre
\(\mathcal{B}(\mathbb{R})\). \(\mathbb{P}_X\) et \(\mathbb{P}_Y\) étant
des mesures finies, la proposition C.3.12 du cours de calcul intégral
nous permet d'assurer que \(\mathbb{P}_X = \mathbb{P}_Y\) pour tous les
éléments de \(\mathcal{B}(\mathbb{R})\), autrement dit que ces deux
mesures de probabilités sont identiques sur
\((\mathbb{R},\mathcal{B}(\mathbb{R}))\).\hfill$\blacksquare$
\end{preuve}

\addcontentsline{toc}{subsection}{Caractérisation de la fonction de
répartition}
\begin{theoreme}[{Caractérisation de la fonction de
répartition}]\phantomsection\label{I-theofdr}

Une fonction \(F\) est la fonction de répartition d'une unique
probabilité \(\mathbb{P}_X\) sur
\((\mathbb{R},\mathcal{B}(\mathbb{R}))\) si et seulement si elle vérifie
les trois conditions suivantes~:

\begin{enumerate}
\def\labelenumi{\arabic{enumi}.}
\tightlist
\item
  elle est croissante,
\item
  elle est continue à droite,
\item
  \(\lim\limits_{x \to -\infty} F(x) = 0, \lim\limits_{x \to +\infty} F(x) = 1\).
\end{enumerate}
\end{theoreme}

\begin{preuve}
La première assertion est immédiate d'après la
\protect\hyperref[I-deffdr]{définition (p.~\pageref*{I-deffdr})}. Pour
la seconde, on remarque que si \(x_n\) décroît vers \(x\), alors
\(]-\infty,x_n]\) décroît vers \(]-\infty,x]\) et donc \(F(x_n)\)
décroît vers \(F(x)\) par le
\protect\hyperref[I-continuitemonotone]{théorème de la continuité
monotone (p.~\pageref*{I-continuitemonotone})}. La troisième assertion
se montre de manière analogue en remarquant que \(]-\infty,x]\) décroît
vers \(\varnothing\) (resp. croît vers \(\mathbb{R}\)) lorsque \(x\)
décroît vers \(-\infty\) (resp. croît vers \(+\infty\)).

Pour la réciproque, on se reportera au chapitre 7 de Jacod and Protter
(\protect\hyperlink{refI-Jacod}{2003}).\hfill$\blacksquare$
\end{preuve}

\addcontentsline{toc}{subsection}{Quelques propriétés}
\begin{remarque}[{Quelques
propriétés}]\phantomsection\label{I-quelques-propriuxe9tuxe9s}

Comme \(F\) est croissante, elle admet une limite à gauche en chaque
point notée \(F(x^-)\). En remarquant que
\(]-\infty,y[\, = \lim\limits_{n \to +\infty}]-\infty,y_n[\) si \(y_n\)
tend vers \(y\) par valeurs décroissantes, on obtient pour \(x < y\)~:

\begin{itemize}
\tightlist
\item
  \(\mathbb{P}_X(]x,y]) = \mathbb{P}( x < X \leq y) = F(y) - F(x)\)
\item
  \(\mathbb{P}_X(]x,y[) = \mathbb{P}( x < X < y) = F(y-) - F(x)\)
\item
  \(\mathbb{P}_X([x,y]) = \mathbb{P}( x \leq X \leq y) = F(y) - F(x^-)\)
\item
  \(\mathbb{P}_X([x,y[) = \mathbb{P}( x \leq X < y) = F(y-) - F(x^-)\)
\end{itemize}

En particulier, \(\mathbb{P}_X(\{x\}) = F(x) - F(x^-)\) est le
\textbf{saut} de la fonction \(F\) au point \(x\). On a donc
\(\mathbb{P}_X(\{x\}) = 0\) pour tout \(x\) si et seulement si \(F\) est
continue en tout point.
\end{remarque}

\addcontentsline{toc}{subsection}{Nécessité des tribus}
\begin{remarque}[{Nécessité des
tribus}]\phantomsection\label{I-nuxe9cessituxe9-des-tribus}

Le \protect\hyperref[I-carac]{théorème
ci-dessus (p.~\pageref*{I-carac})} explique pourquoi, d'un point de vue
strictement mathématique, il est nécessaire d'introduire les tribus en
probabilités, malgré la complexité que cela engendre.

Plus concrètement, considérons l'exemple suivant~: soit
\(\Omega = [0,1]\) et \(\mathbb{P}\) telle que
\(\mathbb{P}(]a,b]) = b-a\) pour \(0\leq a\leq b\leq 1\) (il s'agit de
la loi uniforme sur {[}0,1{]}). C'est une probabilité naturelle qui
assigne à tout intervalle sa longueur comme probabilité. Supposons
maintenant que l'on souhaite étendre de manière unique \(\mathbb{P}\)
aux \(2^{[0,1]}\) éléments de \(\mathcal{P}([0,1])\) de manière à ce que
\(\mathbb{P}(\Omega) =1\) et
\(\mathbb{P}\left(\cup_{n\in\mathbb{N}^\star} A_n\right) = \sum_{n\in\mathbb{N}^\star} \mathbb{P}(A_n)\)
pour toute suite \((A_n)_{n\in\mathbb{N}^\star}\) tels que
\(\mathcal{A}_n\cap A_m = \varnothing\) pour \(n \neq m\). On peut
prouver qu'un tel \(\mathbb{P}\) n'existe pas. \(\mathcal{P}([0,1])\)
est trop ``grand'' pour définir un tel \(\mathbb{P}\). Il contient en
particulier des ensembles non mesurables.

Si l'on voulait travailler avec la tribu
\(\mathcal{A}= \mathcal{P}(\mathbb{R})\), il n'existerait que très peu
de probabilités sur \(\mathbb{R}\), à savoir les probabilités discrètes
que l'on décrit rapidement ci-dessous.
\end{remarque}

%\addcontentsline{toc}{subsection}{Probabilités discrètes}
\begin{exemple}[{Probabilités
discrètes}]\phantomsection\label{I-ex.discret}

\begin{enumerate}
\def\labelenumi{\arabic{enumi}.}
\item
  Les masses de Dirac (ou \textbf{mesures} de Dirac).

  \(X\) est identiquement égale à \(a \in \mathbb{R}\). Alors sa loi est
  la mesure de Dirac en \(a\), la probabilité \(\mathbb{P}_X\) sur
  \(\mathbb{R}\) qui vérifie pour \(A \in \mathcal{P}(\mathbb{R})\)
  \begin{equation*}
      \mathbb{P}_X(A) = \left\{ \begin{array}{ll}
      1  &\text{si } a \in A, \\
      0 &\text{sinon.}
      \end{array}
      \right.
  \end{equation*} Sa fonction de répartition est
  \(F(x) = 1_{[a,+\infty[}(x)\).
\item
  Les probabilités portées par \(\mathbb{N}\).

  \(X\) est à valeurs dans \(\mathbb{N}\). Comme \(\mathbb{N}\) est une
  partie de \(\mathbb{R}\), toute probabilité sur \(\mathbb{N}\) peut
  être considérée comme une probabilité sur \(\mathbb{R}\) qui ne
  ``charge'' que \(\mathbb{N}\). Plus précisément, si \(Q\) est la loi
  de probabilité de \(X\) sur \(\mathbb{N}\), on définit son
  ``extension'' \(\mathbb{P}_X\) à \(\mathbb{R}\) en posant
  \(\mathbb{P}_X(A) = Q(A\cap \mathbb{N})\). Si \(q_n = Q(\{n\})\) pour
  \(n \in \mathbb{N}\), la fonction de répartition \(F\) de
  \(\mathbb{P}_X\) est \begin{equation*}
  F(x) = \left\{ \begin{array}{ll}
  0  &\text{si } x <0, \\
  \sum_{i=0}^{\lfloor x \rfloor} q_i &\text{sinon,}
  \end{array}
  \right.
  \end{equation*} où \(\lfloor \cdot \rfloor\) désigne la partie
  entière. A titre d'exemple, on représente ci-dessous les fonctions de
  répartitions de la loi binomiale et de la loi de Poisson.
  \begin{center}
    \includegraphics{CdfBinom.tex.pdf}
  \end{center}
  \begin{center}
    \includegraphics{CdfPoisson.tex.pdf}
  \end{center}  
    
\item
  Les probabilités discrètes.

  Plus généralement, si \(X\) prend ses valeurs dans une partie finie ou
  dénombrable \(E\) de \(\mathbb{R}\), la loi \(\mathbb{P}_X\) de \(X\)
  est caractérisée pour tout \(i \in E\) par
  \(q_i = \mathbb{P}_X({i}) = \mathbb{P}(X=i)\). La fonction de
  répartition \(F\) de \(X\) est alors
  \[F(x) = \sum_{\substack{i \in E \\ i \leq x}} q_i,\] avec la
  convention qu'une somme ``vide'' vaut 0. On retrouve bien l'exemple 2
  si \(E = \mathbb{N}\). On voit que \(F\) est \textbf{purement
  discontinue} au sens où elle est complètement caractérisée par ses
  sauts \(\triangle F(x) = F(x) - F(x^-)\)~:
  \[F(x) = \sum_{y\leq x} \triangle F(y).\]
\end{enumerate}

Il existe bien d'autres probabilités, non discrètes, sur \(\mathbb{R}\).
Le paragraphe suivant est consacré à un exemple très important, celui
des variables aléatoires dont la loi admet une densité.
\end{exemple}

\begin{figure}
\centering
\includegraphics{CdfBinom.tex.pdf}
\caption{fonction de répartition de la loi
binomiale}\label{I-fig-CdfBinom.tex}
\end{figure}

\begin{figure}
\centering
\includegraphics{CdfPoisson.tex.pdf}
\caption{fonction de répartition de la loi de
Poisson}\label{I-fig-CdfPoisson.tex}
\end{figure}

\section{Variables aléatoires réelles à
densité}\label{I-variables-aluxe9atoires-ruxe9elles-uxe0-densituxe9}

\addcontentsline{toc}{subsection}{Densité de probabilité}
\begin{definition}[{Densité de
probabilité}]\phantomsection\label{I-defdens}

Une fonction réelle \(f\) sur \(\mathbb{R}\) est une \emph{densité de
probabilité} (ou plus simplement une \emph{densité}) si elle est
positive, intégrable et vérifie \[\int_\mathbb{R}f(x)\, dx = 1,\]

où \(dx\) désigne la mesure de Lebesgue sur \(\mathbb{R}\).

Si \(f\) est une densité, la fonction
\[F(x) =\int_{-\infty}^x f(y)\, dy\] est la fonction de répartition
d'une probabilité sur \(\mathbb{R}\).
\end{definition}

\addcontentsline{toc}{subsection}{Variable aléatoire réelle à densité}
\begin{definition}[{Variable aléatoire réelle à
densité}]\phantomsection\label{I-va.densituxe9}

Soit \(X\) une variable aléatoire. On dit que \(X\) a une \emph{loi de
densité \(f\)} (ou par abus de language ``est de densité \(f\)''), si
\(\mathbb{P}_X\) admet la densité \(f\) et donc si pour tout réel \(x\),
\[ \mathbb{P}(X\leq x) = \int_{-\infty}^x f(y) dy.\]
\end{definition}

\addcontentsline{toc}{subsection}{Fonction de répartition d'une variable
aléatoire à densité}
\begin{proposition}[{Fonction de répartition d'une variable aléatoire à
densité}]\phantomsection\label{I-fdr-va-dens}

Soit \(X\) de loi de densité \(f\).

\begin{enumerate}
\def\labelenumi{\arabic{enumi}.}
\tightlist
\item
  Sa fonction de répartition \(F\) est continue, de sorte que
  \[\mathbb{P}(X=x) = 0,~~~\forall x \in \mathbb{R}.\]
\item
  La fonction \(F\) est dérivable en tout point \(x\) où \(f\) est
  continue et \[F'(x) = f(x).\]
\item
  A l'inverse, si la fonction de répartition \(F\) de \(X\) est
  dérivable, ou seulement continue partout et dérivable par morceaux,
  alors \(X\) admet la densité \[F'(x) = f(x).\]
\end{enumerate}
\end{proposition}

Il existe bien sûr des variables aléatoires qui n'ont pas de densité~:
c'est le cas des variables discrètes données en exemple
\protect\hypertarget{I-ex.discret}{}{ci-dessus}. Il existe aussi des cas
``mixtes''~: soient d'une part \(f\) une fonction positive, intégrable
et d'intégrale strictement positive et d'autre part une partie finie ou
dénombrable \(E\) de \(\mathbb{R}\) et des poids \(p_i>0\) indexés par
\(i \in E\), tels que~:
\[ \int_\mathbb{R}f(x)\, dx + \sum_{i\in E} p_i = 1.\] Alors la fonction
\[ F(x) = \int_{-\infty}^x f(x)\, dx + \sum_{\substack{i\in E \\ i \leq x}} p_i\]
est une fonction de répartition, et la probabilité associée
\(\mathbb{P}_X\) n'admet pas de densité et n'est pas non plus discrète.

La mesure de probabilité associée s'écrit
\(\mathbb{P}_X = f\lambda + \sum_{i\in E} p_i \delta_i\) où \(\lambda\)
est la mesure de Lebesgue.

\begin{remarque}\phantomsection\label{I-remarques-1}
\begin{enumerate}
\def\labelenumi{\arabic{enumi}.}
\item
  La fonction de répartition est entièrement déterminée par la
  probabilité \(\mathbb{P}\). Il n'en est pas de même de la densité
  lorsqu'elle existe~: si en effet on a
  \(F(x) =\int_{-\infty}^x f(y)\, dy\) et si on pose \(g(x) = f(x)\) si
  \(x \notin E\) et \(g(x) =f(x)+1\) (par exemple) si \(x\in E\), où
  \(E\) est un ensemble négligeable, alors \(g\) est encore une densité
  de \(\mathbb{P}\). Autrement dit, la densité
\item
  Une interprétation intuitive de la densité \(f\) de \(\mathbb{P}\)~:
  si \(dx\) est un petit accroissement de la variable \(x\), on a (si du
  moins \(f\) est continue en \(x\))~:
  \[ f(x) \sim \frac{\mathbb{P}([x,x+dx])}{dx}.\]
\end{enumerate}
\end{remarque}

\runinhead{Exercice -- Durée de
fonctionnement (\(\mathord{\bullet}\))}\phantomsection\addcontentsline{toc}{subsubsection}{Durée
de fonctionnement}\label{I-ex.expo}
On suppose que la durée de fonctionnement, en heures, d'un ordinateur
avant sa première panne est une variable aléatoire positive de loi
exponentielle de paramètre 1/100, de densité donnée par

\[f(x) = \left\{ \begin{array}{ll}
        \frac{1}{100}\exp\left(-\frac{x}{100}\right) & x \geq 0, \\
        0 & x < 0.
        \end{array}
        \right.\]

\begin{enumerate}
\def\labelenumi{\arabic{enumi}.}
\tightlist
\item
  Calculer la probabilité que cette durée de fonctionnement \(X\) soit
  comprise entre 50 et 150 heures.
\item
  Calculer la probabilité que l'ordinateur fonctionne moins de 100
  heures.
\item
  Calculer la probabilité que l'ordinateur fonctionne exactement 100
  heures avant sa première panne.
\end{enumerate}

\solref{I-answer-ex.expo}

\runinhead{Exercice -- Précipitation (\(\mathord{\bullet}\))}\phantomsection\addcontentsline{toc}{subsubsection}{Précipitation}\label{I-ex.pluie}
On souhaite modéliser la hauteur de précipitation (en mm) tombée sur
Paris en une journée par une variable aléatoire \(X\). On considère que
cette hauteur suit une loi exponentielle de paramètre \(1/10\) lorsqu'il
pleut et qu'il pleut en moyenne un jour sur deux.

\begin{enumerate}
\def\labelenumi{\arabic{enumi}.}
\tightlist
\item
  Donner la Fonction de répartition de \(X\).
\item
  \(X\) admet-elle une densité\,?
\end{enumerate}

\solref{I-answer-ex.pluie}

\addcontentsline{toc}{subsection}{Quelques lois à densités}
\begin{exemple}[{Quelques lois à
densités}]\phantomsection\label{I-exampledens}

\begin{enumerate}
\def\labelenumi{\arabic{enumi}.}
\item
  La \emph{loi uniforme} sur \([a,b]\), avec \(a,b \in \mathbb{R}\) tels
  que \(a < b\), notée \(\mathcal{U}_{[a,b]}\), est la probabilité
  \(\mathbb{P}\) qui admet la densité \[f(x) = \left\{\begin{array}{ll}
          \frac{1}{b-a} & \text{si } a \leq x \leq b,\\
          0 & \text{sinon}.
          \end{array}\right.\] 
          
          \begin{center}
          \includegraphics{PdfUnif.tex.pdf}
          \end{center}
  En
  vertu de la remarque 1. ci-dessus, on aurait pu choisir
  \(f(a) = f(b) = 0\). Au vu de l'interprétation 2. ci-dessus, le fait
  que \(f\) soit constante sur \([a, b]\) exprime que si on choisit un
  point selon la probabilité uniforme, on a ``autant de chances'' de
  tomber au voisinage de chaque point de l'intervalle \([a, b]\). Cela
  explique le nom ``uniforme''. On remarque aussi que
  \(\mathbb{P}(\{x\}) = 0\) pour tout \(x\) (comme pour toutes les
  probabilités à densité). On a donc une probabilité nulle de tomber
  exactement en un point \(x\) fixé à l'avance. On vérifie aisément que
  sa fonction de répartition est \[F(x) = \left\{\begin{array}{ll}
          0 & \text{si } x < a,\\
          \frac{x-a}{b-a} & \text{si } a \leq x \leq b,\\
          1 & \text{si } x > b.
          \end{array}\right.\] \begin{center}
          \includegraphics{CdfUnif.tex.pdf}
          \end{center}
\item
  La \emph{loi exponentielle} de paramètre \(\theta\), notée
  \(\mathcal{E}_\theta\), est la loi de densité
  \[f(x) = \left\{\begin{array}{ll}
          0 & \text{si } x < 0,\\
          \theta e^{-\theta x} & \text{sinon}.
          \end{array}\right.\] 
          \begin{center}
          \includegraphics{PdfExp.tex.pdf}
          \end{center}
          et de
  fonction de répartition \[F(x) = \left\{\begin{array}{ll}
          0 & \text{si } x < 0,\\
          1 - e^{-\theta x} & \text{sinon}.
          \end{array}\right.\] \begin{center}
          \includegraphics{CdfExp.tex.pdf}
          \end{center} Dans la
  pratique, on utilise fréquemment la loi exponentielle pour modéliser
  une durée de vie ou le temps d'attente avant l'arrivée d'un événement
  spécifique. Par exemple la durée de vie d'une bactérie, d'un composant
  électronique, la durée d'une conversation téléphonique ou le temps qui
  nous sépare du prochain tremblement de terre peuvent être considérées
  comme des variables aléatoires de loi exponentielle.
\item
  La \emph{loi normale} (ou de \emph{Gauss} ou encore \emph{gaussienne})
  de paramètres \(\mu\) et \(\sigma^2\), notée
  \(\mathcal{N}(\mu,\sigma^2)\), est la loi de densité

  \[f(x) = \frac{1}{\sqrt{2\pi}\sigma}\exp\left(-\frac{(x-\mu)^2}{2\sigma^2}\right)\]
  \begin{center}
  \includegraphics{PdfGauss.tex.pdf}
  \end{center}
   Lorsque \(\mu = 0\) et
  \(\sigma^2 = 1\), on l'appelle la loi normale \emph{centrée réduite}
  pour des raisons que nous expliciterons au prochain chapitre.

  La distribution normale fut introduite par De Moivre en 1733. Celui-ci
  l'utilisa pour approximer une variable aléatoire binomiale quand le
  paramètre \(n\) de celle-ci était grand. Ce résultat fut ensuite
  progressivement généralisé par Laplace et d'autres confrères pour
  devenir le théorème connu sous le nom de théorème de la limite
  centrale, qui sera démontré au Chapitre 4. Ce théorème est l'un des
  plus importants de la théorie des probabilités et prouve que de très
  nombreux phénomènes aléatoires suivent approximativement une loi
  normale. On peut citer à titre d'exemple la taille d'un individu
  choisi au hasard, les composantes de la vitesse d'une molécule de gaz
  ou l'erreur de mesure d'une quantité physique.

  Il est difficile de faire des calculs avec la loi normale car sa
  densité n'admet pas de primitive explicite. Aussi des tables
  numériques ont-elles été construites pour permettre aux utilisateurs
  d'obtenir très rapidement des valeurs numériques. Elles sont
  disponibles dans la plupart des logiciels et permettent entre autre de
  représenter sa fonction de répartition.

\begin{center}
\includegraphics{CdfGauss.tex.pdf}
\captionof{figure}{fonction de répartition de la loi normale}
\end{center} 
\end{enumerate}

Nous aurons l'occasion de voir par la suite un grand nombre d'autres
exemples de probabilités à densité.
\end{exemple}

\begin{figure}
\centering
\includegraphics{PdfUnif.tex.pdf}
\caption{Densité de probabilité de la loi
uniforme}\label{I-fig-PdfUnif.tex}
\end{figure}

\begin{figure}
\centering
\includegraphics{CdfUnif.tex.pdf}
\caption{fonction de répartition de la loi
uniforme}\label{I-fig-CdfUnif.tex}
\end{figure}

\begin{figure}
\centering
\includegraphics{PdfExp.tex.pdf}
\caption{Densité de probabilité de la loi
exponentielle}\label{I-fig-PdfExp.tex}
\end{figure}

\begin{figure}
\centering
\includegraphics{CdfExp.tex.pdf}
\caption{fonction de répartition de la loi
exponentielle}\label{I-fig-CdfExp.tex}
\end{figure}

\begin{figure}
\centering
\includegraphics{PdfGauss.tex.pdf}
\caption{Densité de probabilité de la loi
normale}\label{I-fig-PdfGauss.tex}
\end{figure}

\begin{figure}
\centering
\includegraphics{CdfGauss.tex.pdf}
\caption{fonction de répartition de la loi
normale}\label{I-fig-CdfGauss.tex}
\end{figure}

\chapter{Exercices
complémentaires}\label{I-exercices-compluxe9mentaires}

\section{Indépendance et
conditionnement}\label{I-induxe9pendance-et-conditionnement}

\runinhead{Question
1}\phantomsection\addcontentsline{toc}{subsubsection}{Question
1}\label{I-ic-1}
Dans quelles circonstances particulières deux événements \(A\) et \(B\)
tels que \(A\cup B\) soit presque sûr sont-ils
indépendants\,?\solref{I-answer-ic-1}

\runinhead{Question
2}\phantomsection\addcontentsline{toc}{subsubsection}{Question
2}\label{I-ic-2}
Soit \(A\), \(B\) et \(C\) trois événements indépendants \textbf{2 à 2},
avec \(\mathbb{P}(C)>0\). Donner une condition nécessaire et suffisante
sur \(A\), \(B\) et \(C\) pour que \(A\) et \(B\) soient indépendants
relativement à la probabilité conditionnelle \(\mathbb{P}(\,\,|C)\), à
la place de \(\mathbb{P}\).\solref{I-answer-ic-2}

\runinhead{Question
3}\phantomsection\addcontentsline{toc}{subsubsection}{Question
3}\label{I-ic-3}
Etablir que deux événements \(A\) et \(B\) sont indépendants ssi
\[ \mathbb{P}(A\cap B)\mathbb{P}(A^c \cap B^c) = \mathbb{P}(A \cap B^c)\mathbb{P}(A^c\cap B)\]

\solref{I-answer-ic-3}

\section{Problème de Monty Hall}\label{I-mh}

Dans un jeu télévisé, vous êtes confrontés au problème suivant~: devant
vous se trouvent 3 portes fermées. Derrière l'une d'entre elle se trouve
une voiture, derrière chacune des deux autres, une chèvre. Vous désignez
au hasard l'une d'entre elle. Le présentateur, qui sait où se trouve la
voiture, ouvre alors une porte, qui n'est ni celle que vous avez
choisie, ni celle qui cache la voiture. Il vous offre alors la
possibilité de réviser votre choix. Que choisissez-vous\,?

Calculer la probabilité de remporter la voiture selon les deux
stratégies (changer ou non son choix de porte).\solref{I-answer-mh}

\section{Mesure de Probabilité}\label{I-mesure-de-probabilituxe9}

Soit \((\Omega,\mathcal{A},\mathbb{P})\) un espace probabilisé et \(X\)
une variable aléatoire réelle définie sur
\((\Omega,\mathcal{A},\mathbb{P})\) telle que \(X \geq 0\) p.s. et
\(\mathbb{E}(X) = \int_\Omega X(\omega) d\mathbb{P}(\omega) = 1\). (on
reviendra sur cette définition dans le chapitre suivant).

On définit \(Q : \mathcal{A}\to \mathbb{R}\) par
\(Q(A) = \mathbb{E}(X\,1_A) = \int_{\Omega} X(\omega) 1_A(\omega)d\mathbb{P}(\omega)\)
pour \(A\in \mathcal{A}\).

\runinhead{Question
1}\phantomsection\addcontentsline{toc}{subsubsection}{Question
1}\label{I-mes-proba1}
Montrer que \(Q\) est une mesure de probabilité sur
\((\Omega,\mathcal{A})\).\solref{I-answer-mes-proba1}

\runinhead{Question
2}\phantomsection\addcontentsline{toc}{subsubsection}{Question
2}\label{I-mes-proba2}
Montrer que si \(\mathbb{P}(A) = 0\), alors \(Q(A) = 0\). Donner un
exemple montrant que \(Q(A) = 0\) n'implique pas en général que
\(\mathbb{P}(A) = 0\).\solref{I-answer-mes-proba2}

\runinhead{Question
3}\phantomsection\addcontentsline{toc}{subsubsection}{Question
3}\label{I-mes-proba3}
En notant \(\mathbb{E}_Q\), l'espérance par rapport à \(Q\), montrer que
\(\mathbb{E}_Q(Y) = \int_{\Omega} Y(\omega) Q(d\omega) = \mathbb{E}_\mathbb{P}(XY)\)
pour toute variable aléatoire réelle \(Y\) définie sur
\((\Omega,\mathcal{A},\mathbb{P})\) positive ou
bornée.\solref{I-answer-mes-proba3}

\runinhead{Question
4}\phantomsection\addcontentsline{toc}{subsubsection}{Question
4}\label{I-mes-proba4}
Supposons que \(\mathbb{P}(X > 0) = 1\).

\begin{enumerate}
\def\labelenumi{\roman{enumi})}
\item
  Montrer que \(\frac{1}{X}\) est intégrable pour \(Q\).
\item
  Soit \(R : \mathcal{A}\to \mathbb{R}\) telle que
  \(\forall A \in \mathcal{A}\,\, R(A) = \mathbb{E}_Q(\frac{1}{X}1_A)\).
  Montrer que \(R\) est exactement la probabilité \(\mathbb{P}\).
\item
  Montrer que \(Q(A) = 0\) implique \(\mathbb{P}(A) = 0\).
\end{enumerate}

\solref{I-answer-mes-proba4}

\section{Queues de distributions}\label{I-queues-de-distributions}

Soit une densité \(f\) sur \(\mathbb{R}\) telle que la fonction
\(h: x \in \mathbb{R}\mapsto x\,f(x)\) est intégrable. Cette hypothèse
sera appelée \((\mathcal{H})\). On note \(F\) la fonction de répartition
correspondante. Nous allons montrer que\footnote{Soient
  \(a \in [-\infty,+\infty]\) et \(f\) et \(g\) deux fonctions définies
  sur un même voisinage \(I\) de \(a\). On dit que \(f\) est
  \emph{négligeable} par rapport à \(g\) (ou que \(g\) est
  \emph{prépondérante} par rapport à \(f\)) au voisinage de \(a\) et on
  note note \(f(x) \underset{x\to a}{=} o(g(x))\) s'il existe une
  fonction \(\epsilon\) définie sur \(I\) telle que
  \(\epsilon(x) \xrightarrow[x\to a]{} 0\) et \(f = \epsilon \times g\).}
\[\left|\begin{array}{rl} 1-F(x) \underset{x\to+\infty}{=} o\left(\dfrac{1}{x}\right), & (1)\\[1em] F(x) \underset{x\to-\infty}{=} o\left(\dfrac{1}{x}\right).& (2) \end{array} \right. \]

\runinhead{Etude en
\(+\infty\)}\phantomsection\addcontentsline{toc}{subsubsection}{Etude en
\(+\infty\)}\label{I-dlfdr-pi}

\begin{enumerate}
\def\labelenumi{\arabic{enumi}.}
\item
  Vérifier que pour tout \(u > 0\),
  \[\int_{-\infty}^{+\infty} h(x)\,dx \geq \int_{-\infty}^u h(x)\,dx + u\,\left(1-F(u)\right).\]
\item
  En déduire \((1)\).
\end{enumerate}

\solref{I-answer-dlfdr-pi}

\runinhead{Etude en
\(-\infty\)}\phantomsection\addcontentsline{toc}{subsubsection}{Etude en
\(-\infty\)}\label{I-dlfdr-mi}

\begin{enumerate}
\def\labelenumi{\arabic{enumi}.}
\setcounter{enumi}{2}
\tightlist
\item
  En vous inspirant du raisonnement précédent, montrer \((2)\).
\end{enumerate}

\solref{I-answer-dlfdr-mi}

\runinhead{Interprétation}\phantomsection\addcontentsline{toc}{subsubsection}{Interprétation}\label{I-dlfdr-interpret}

\begin{enumerate}
\def\labelenumi{\arabic{enumi}.}
\setcounter{enumi}{3}
\tightlist
\item
  Proposez une interprétation de ce résultat.
\end{enumerate}

\solref{I-answer-dlfdr-interpret}

\runinhead{Etude de la
réciproque}\phantomsection\addcontentsline{toc}{subsubsection}{Etude de
la réciproque}\label{I-dlfdr-reciproque}
Considérons maintenant la fonction de répartition
\[ F: x \in \mathbb{R}\mapsto \left|\begin{array}{ll} 1 - \dfrac{e}{x\,\ln(x)} & \text{si } x \geq e,\\ 0 & \text{si } x <e, \end{array}\right. \]

\begin{enumerate}
\def\labelenumi{\arabic{enumi}.}
\setcounter{enumi}{4}
\item
  Vérifier que cette fonction possède bien une densité \(f\) et
  l'expliciter.
\item
  A-t-on les développements asymptotiques \((1)\) et \((2)\)\,?
\item
  La fonction \(h\) correspondante est-elle absolument intégrable\,?
\item
  Qu'en concluez-vous\,?
\end{enumerate}

\solref{I-answer-dlfdr-reciproque}

\runinhead{Liens entre les développements
asymptotiques}\phantomsection\addcontentsline{toc}{subsubsection}{Liens
entre les développements asymptotiques}\label{I-dlfdr-liensdl}
Soient les fonctions de répartition
\[F: x\in\mathbb{R}\mapsto \left|\begin{array}{ll} 1 - \dfrac{1}{\ln(x)} & \text{si } x \geq e,\\ 0 & \text{sinon,}\end{array}\right.\]
et \(G : x\in\mathbb{R}\mapsto 1 - F(-x)\).

\begin{enumerate}
\def\labelenumi{\arabic{enumi}.}
\setcounter{enumi}{8}
\tightlist
\item
  Que peut-on dire sur les éventuels liens entre \((1)\) et \((2)\) à
  l'aide de ces deux fonctions\,?
\end{enumerate}

\solref{I-answer-dlfdr-liensdl}

\section{Densité et fonction de répartition d'une loi
Normale}\label{I-densituxe9-et-fonction-de-ruxe9partition-dune-loi-normale}

On considère la densité d'une loi Normale centrée réduite~:
\[f : x\in \mathbb{R}\mapsto \dfrac{1}{\sqrt{2\,\pi}}\,\exp\left\{-\dfrac{x^2}{2}\right\},\]
où \(f(-\infty) := \lim\limits_{x\to-\infty} f(x)\) et
\(f(+\infty) := \lim\limits_{x\to+\infty} f(x)\).

\runinhead{Propriétés
générales}\phantomsection\addcontentsline{toc}{subsubsection}{Propriétés
générales}\label{I-fdrgaussprop}

\begin{enumerate}
\def\labelenumi{\arabic{enumi}.}
\item
  Faire l'étude de \(f\) (domaine, parité, limites aux bornes,
  dérivabilité, variations, convexité/concavité).
\item
  Donner la définition de la fonction de répartition \(F\) associée à
  \(f\). En donner une expression faisant apparaître la \emph{fonction
  d'erreur} (qui est une fonction spéciale)
  \[\text{erf} : x\in\mathbb{R}\mapsto \dfrac{2}{\sqrt{\pi}}\,\int_{0}^x e^{-u^2}\,du.\]
\end{enumerate}

\solref{I-answer-fdrgaussprop}

\runinhead{Encadrement de \(1 - F(x)\) pour tout
\(x > 0\)}\phantomsection\addcontentsline{toc}{subsubsection}{Encadrement
de \(1 - F(x)\) pour tout \(x > 0\)}\label{I-fdrgaussenca}
Nous allons maintenant démontrer la propriété suivante~:
\(\forall\,x > 0\)
\[ \dfrac{x}{1+x^2}\,f(x) \leq 1 - F(x) \leq \dfrac{f(x)}{x}. \]

\begin{enumerate}
\def\labelenumi{\arabic{enumi}.}
\setcounter{enumi}{2}
\item
  En observant que \(u\geq x >0\) implique \(\dfrac{u}{x} \geq 1\),
  montrer l'inégalité de droite.
\item
  En observant que \(u\geq x > 0\) implique \(u^{-2} \leq x^{-2}\),
  montrer que \(\forall\,x>0\),
  \[\left(1+\dfrac{1}{x^2}\right)\,\left(1-F(x)\right) \geq \dfrac{f(x)}{x}.\]
  En déduire l'inégalité de gauche.
\end{enumerate}

\solref{I-answer-fdrgaussenca}

\runinhead{Equivalent de \(1-F(x)\) quand
\(x\to+\infty\).}\phantomsection\addcontentsline{toc}{subsubsection}{Equivalent
de \(1-F(x)\) quand \(x\to+\infty\).}\label{I-fdrgaussequi}

\begin{enumerate}
\def\labelenumi{\arabic{enumi}.}
\setcounter{enumi}{4}
\tightlist
\item
  Déduire de l'encadrement prédédent un équivalent de \(1-F(x)\) lorsque
  \(x\to+\infty\).
\end{enumerate}

\solref{I-answer-fdrgaussequi}

\begin{enumerate}
\def\labelenumi{\arabic{enumi}.}
\setcounter{enumi}{5}
\tightlist
\item
  Comparer ce résultat à celui de l'exercice sur les développements
  limités pour les fonctions de répartition. En particulier,
  \(h:x\in[-\infty,+\infty] \mapsto x\,f(x)\) est-elle intégrable\,?
\end{enumerate}

\runinhead{Loi Normale
générale.}\phantomsection\addcontentsline{toc}{subsubsection}{Loi
Normale générale.}\label{I-fdrgaussgen}
Considérons maintenant la densité plus générale
\[g(x) = \dfrac{1}{\sqrt{2\,\pi}\,\sigma}\, \exp\left\{-\dfrac{(x-\mu)^2}{2\,\sigma^2} \right\},\]
où \(\mu \in \mathbb{R}\) et \(\sigma^2 > 0\). On note \(G\) sa fonction
de répartition associée.

\begin{enumerate}
\def\labelenumi{\arabic{enumi}.}
\setcounter{enumi}{6}
\tightlist
\item
  En réécrivant \(g\) en fonction de \(f\), déduire des questions
  précédentes un équivalent de \(1-G(x)\) lorsque \(x\to+\infty\).
\end{enumerate}

\solref{I-answer-fdrgaussgen}

\section{Loi de vie et de mort}\label{I-loi-de-vie-et-de-mort}

La durée de vie d'un être vivant ou d'un matériel peut-être assimilée à
une variable aléatoire strictement positive \(T\). Dans ce cadre, on
peut définir les notions suivantes~:

\begin{itemize}
\tightlist
\item
  Loi de vie a priori~: il s'agit de la loi du temps \(T\) caractérisée
  par fonction de répartition complémentaire \[G(t) = \mathbb{P}(T>t)\]
\item
  Loi de survie après \(t_0 \geq 0\)~: il s'agit de la loi du temps
  \(T-t_0\) qu'il lui reste à vivre, sachant qu'il était encore en vie à
  \(t_0\), de fonction de répartition complémentaire
  \[G_{t_0}(t) =\mathbb{P}(T>t+t_0|T>t_0)\]
\end{itemize}

On dira que la loi de vie satisfait la propriété de non vieillissement
(ou d'absence de mémoire) si la loi de survie et la loi de vie sont
égales~:
\[\forall t \geq 0 \text{ et } \forall t_0 \geq 0, G_{t_0}(t) = G(t)\]

\runinhead{Question
1}\phantomsection\addcontentsline{toc}{subsubsection}{Question
1}\label{I-viemort1}
Exprimer la loi de survie à partir de la loi de vie a
priori.\solref{I-answer-viemort1}

\runinhead{Question
2}\phantomsection\addcontentsline{toc}{subsubsection}{Question
2}\label{I-viemort2}
Montrer qu'une variable aléatoire \(T\) satisfait la propriété de
non-vieillissement si et seulement si elle est de loi
exponentielle.\solref{I-answer-viemort2}

\runinhead{Question
3}\phantomsection\addcontentsline{toc}{subsubsection}{Question
3}\label{I-viemort3}
On suppose que \(T\) admet une densité continue sur
\(\mathbb{R}_+^\ast\). Montrer que l'on peut définir le taux de mort à
l'instant \(t\) par
\[D(t) = \lim_{\Delta t \to 0} \frac{1}{\Delta t}\mathbb{P}(t < T \leq t + \Delta t | T >t)\]
et exprimer \(G\) en fonction de \(D\).

Quelles lois correspondent-elles à \(D\) constant\,?
\solref{I-answer-viemort3}

\begin{center}\rule{0.5\linewidth}{0.5pt}\end{center}



\chapter{Références}\label{I-ruxe9fuxe9rences}

\hypertarget{refs}{}
\begin{cslreferences}
\leavevmode\hypertarget{refI-polygarnier}{}%
Garnier, Josselin. 2018. \emph{Aléatoire}.
\url{http://josselin-garnier.org/wp-content/uploads/2018/03/MAP311_2018.pdf}.

\leavevmode\hypertarget{refI-Jacod}{}%
Jacod, J., and P. Protter. 2003. \emph{L'essentiel En Théorie Des
Probabilités}. Cassini.
\url{https://hal.archives-ouvertes.fr/hal-00104956}.

\leavevmode\hypertarget{refI-polyponts}{}%
Jourdain, Benjamin. 2016. \emph{Probabilités et Statistique}. Ellipses.
\url{http://cermics.enpc.fr/~jourdain/probastat/poly.pdf}.

\leavevmode\hypertarget{refI-matheron}{}%
Matheron, Georges. 1978. \emph{Estimer et choisir: essai sur la pratique
des probabilités}. Fontainebleau, France: Ecole nationale supérieure des
mines de Paris.
\url{http://cg.ensmp.fr/bibliotheque/public/MATHERON_Ouvrage_00208.pdf}.

\leavevmode\hypertarget{refI-tebaldi}{}%
Tebaldi, Claudia, and Reto Knutti. 2007. ``The Use of the Multi-Model
Ensemble in Probabilistic Climate Projections.'' \emph{Philosophical
Transactions of the Royal Society A: Mathematical, Physical and
Engineering Sciences} 365 (1857): 2053--75.
\url{https://doi.org/10.1098/rsta.2007.2076}.
\end{cslreferences}
